\documentclass[11pt]{article}
\usepackage[margin=1in]{geometry}
\usepackage{amsmath,amssymb,amsthm,mathtools}
\IfFileExists{lmodern.sty}{\usepackage{lmodern}}{}
\usepackage[expansion=false]{microtype}
\usepackage{booktabs,array}
\usepackage{enumitem}
\usepackage{xcolor}
\usepackage[colorlinks=true,linkcolor=blue!50!black,citecolor=green!35!black,urlcolor=blue!50!black]{hyperref}
\usepackage[nameinlink]{cleveref}

\newtheorem{theorem}{Theorem}[section]
\newtheorem{proposition}[theorem]{Proposition}
\newtheorem{lemma}[theorem]{Lemma}
\newtheorem{corollary}[theorem]{Corollary}
\theoremstyle{definition}
\newtheorem{definition}[theorem]{Definition}
\newtheorem{example}[theorem]{Example}
\newtheorem{question}[theorem]{Question}
\theoremstyle{remark}
\newtheorem{remark}[theorem]{Remark}
\newtheorem*{mainthm}{Main results}

\newcommand{\D}{\mathsf{D}}
\newcommand{\R}{\mathbb{R}}
\newcommand{\Z}{\mathbb{Z}}
\newcommand{\Cpl}{\operatorname{Cpl}}
\newcommand{\diag}{\operatorname{diag}}
\newcommand{\T}{\mathsf T}
\newcommand{\C}{\mathsf C}
\newcommand{\Sh}{\mathrm{Sh}}

\title{Coupling Algebras:\\
Gluing, the Square-Zero Cone and the Ring of Contingency Tables}
\author{Vinay K. Awasthi\\
Department of Applied Mathematics and Statistics\\
Johns Hopkins University\thanks{Every statement is checked by \texttt{check\_coupling\_algebra.py} (14 checks, exact
arithmetic except where marked).}}
\date{October 6, 2026, draft 3}

\begin{document}
\maketitle

\begin{abstract}
The algebra that carries the geometry of couplings of finite distributions is the \emph{ring of contingency tables}:
the Ehrhart ring of a coupling polytope with rational margins. It is generated in degree one, its Hilbert function counts
integer tables with scaled margins, and its $\operatorname{Proj}$ is the toric variety of the coupling polytope. Its degree-one
basis is the sample space of Fisher's exact test, and the basic moves that connect that sample space are integer
circulations. For $2\times m$ tables we give the number of tables in closed form on every chamber of margins: it is a
polynomial on each chamber, and crossing a wall changes it by a single binomial term. When the wall cuts a corner of
lattice depth $d$ off the polygon, as in the blow-ups of the toric picture, the count drops by $\binom{d+1}2$, which is the
Riemann--Roch jump of the line bundle. The ring is classical in algebraic statistics; what is new is its place here, as the
object that joins the toric geometry of couplings to exact conditional inference.

The other algebra in the programme, the \emph{gluing algebra} $A_f$ of matrices whose row and column sums are multiples of
$f$, is much smaller than it looks, and this corrects earlier expectations about it. We show $A_f\cong\R\times M_{n-1}(\R)$, with
the product coupling and its complement as central idempotents. Categorically, $A_f$ is the linear span of the monoid of
$\D$-algebra endomorphisms of $\D(Y)$ fixing $f$, with the product coupling the endomorphism through the zero object. The
$A_\infty$-functor of whiskering into the square-zero cone of the circulations exists and is unique for every mass-preserving
linear map, but the cone is quasi-isomorphic to $\R$, and up to homotopy the functor is the mass. Couplings with kernels as
vertical arrows do not form a double category, although squares compose when the middle map is a bijection or merges only
interchangeable values. The tower's context comonad is acted on by the monoid of $f$-fixing kernels, but not linearly, so it does
not see $A_f$.
\end{abstract}

\tableofcontents

%=====================================================================
\section{Introduction}\label{sec:intro}
%=====================================================================

The theory paper \cite{Conditionals} linearizes the coupling 2-cells of the Kleisli category of the finite distribution
monad. For a distribution $f$ on a finite set, it forms the algebra $A_f$ of matrices whose row and column sums are
multiples of $f$, with the gluing product, and the ideal $I_f$ of circulations. Whiskering by a kernel is then a linear
map that is not multiplicative, and its defect is the second Boolean cumulant. Into the square-zero cone of $I_f$,
whiskering is an $A_\infty$-functor with exactly two components \cite[Thm.~5.2]{Conditionals}, also across objects
\cite[Thm.~10.2]{Conditionals}. Two questions were left: why this cone, and whether a larger cone carries more
\cite[Question~9.4]{Conditionals}. A third question, from the attribution paper \cite[Open problem~5]{Attribution}, asks
for a strict structure on the witness layer.

This note has a positive result and a set of corrections. The positive result is that the algebra which carries the coupling
geometry is not $A_f$ but a graded commutative one: the ring of lattice points of the coupling polytope, the ring of
contingency tables. It joins the toric varieties of \cite{Conditionals,Family} to exact conditional inference
\cite{DiaconisSturmfels1998}, and its Hilbert function crosses the walls of margins in the same way that the toric surfaces
are blown up (\cref{prop:twobym}). The ring itself is classical; its role as the meeting point of the two sides is what is new.

The corrections concern the three questions. The cone is homotopy-initial among targets that kill the circulations, but only
because it is equivalent to $\R$: the gluing algebra splits off its circulations as a direct factor. The functor into it is
forced, and up to homotopy it remembers only the mass. And couplings with kernels do not form a double category.

\paragraph{For readers of the theory paper.} The answer to ``why this cone?'' is a correction of the strong form of the
question, not a confirmation of it. The cone of \cite[Thm.~5.2]{Conditionals} is the right construction, and nothing about it
is a choice: by \cref{thm:forced}, every mass-preserving linear map, whiskering included, becomes an $A_\infty$-functor into it in
exactly one way, with $F_2=\varepsilon b_2$ and no higher components. But the universality one might hope for, that the cone is
the initial recipient of the coupling information, is true only in a vacuous form. By \cref{thm:homotopy} the cone is
quasi-isomorphic to $\R$, and by \cref{cor:blind} every homotopy invariant of such a functor is a function of masses. So the
coupling information of whiskering lives at chain level, in $F_2=\varepsilon b_2$, and \cite[Question~9.4(i)]{Conditionals}
has a negative answer for homotopy invariants. Whether a free cone carries chain-level information in the higher Boolean
cumulants remains open.

\begin{mainthm}
\leavevmode
\begin{enumerate}[label=(\Alph*),leftmargin=*]
\item \emph{Structure} (\cref{thm:structure,prop:split}). $A_f\cong\R\times M_{n-1}(\R)$. The product coupling
  $P_f=f\otimes f$ and $e_f=\diag(f)-P_f$ are complementary central idempotents, $I_f=e_fA_f$, and the mass is the
  projection onto $\R P_f$. Across objects the linear category of couplings splits as the indiscrete category on $\R$
  times the category of circulations.
\item \emph{Monad and comonad} (\cref{prop:kleisli,prop:comonad}). Couplings of $f$ with itself are the $\D$-algebra
  endomorphisms of $\D(Y)$ fixing $f$, gluing is Kleisli composition, and $A_f$ is their linear span. The product coupling is
  the endomorphism through the zero object $(1,\delta)$, and the mass is the map to it. The tower's context comonad, whose
  distributive law is the $\D$-algebra structure of the context algebra, is acted on by this monoid, but not affinely at level
  two, so the action does not extend to $A_f$.
\item \emph{Forced functor} (\cref{thm:forced}). For every linear map $F_1$ that preserves mass, there is exactly one
  $A_\infty$-functor into the square-zero cone with first component $F_1$: its second component is the multiplicativity
  defect, and all higher components vanish.
\item \emph{Homotopy content} (\cref{thm:homotopy,cor:blind}). The cone is quasi-isomorphic to $\R$, and the functor composed
  with this quasi-isomorphism is the mass, a strict algebra map. Since $A_f$ is separable, every $A_\infty$-functor out of it
  is homotopic to a strict one, and if it makes $e_f$ a boundary, its homotopy class factors through the mass. So no
  homotopy invariant of such a functor sees contextuality.
\item \emph{No double category} (\cref{prop:double}). With couplings as horizontal arrows, kernels as vertical arrows and
  squares given by pushforward, squares compose horizontally when the middle kernel is a bijection, but not for some
  constant kernels or some non-injective functions. The defect is a circulation.
\item \emph{The ring of contingency tables} (\cref{thm:ehrhart,prop:twobym}). For margins with common denominator $N$, the polytope
  $N\cdot\Cpl(f,g)$ is a lattice polytope with the integer decomposition property. Its Ehrhart ring is generated in degree
  one, its Hilbert function counts integer tables, and its $\operatorname{Proj}$ is the toric variety of the coupling
  polytope. Its degree-one basis is the sample space of Fisher's exact test, connected by basic moves that are integer
  circulations. For $2\times m$ tables the count is a polynomial on each chamber of margins, and across a wall that cuts a
  corner of lattice depth $d$ it drops by $\binom{d+1}2$, the Riemann--Roch jump of a toric blow-up.
\end{enumerate}
\end{mainthm}

\paragraph{What is not claimed.} Earlier drafts of this programme proposed several stronger statements, and each is false.
The square-zero cone is homotopy-initial among targets killing $I_f$, but it is equivalent to $\R$, so initiality says nothing
beyond the mass. Its minimal model, in the sense of $A_\infty$-algebras \cite{Keller2001}, is $\R$ itself. And the witness
layer with kernels as vertical maps is not a double category, strict or lax. The gluing algebra is also not a coordinate ring,
since it is not commutative, and $A_f/I_f\cong\R$ is not the cohomology ring of the toric manifold, whose dimension is the
number of vertices of the polytope. The monadic reading also has limits, listed in \cref{rem:monad-not}: the barycenter is not
the mass projection, there is no ideal filtration matching the precision ladder, and the multiplicativity defect is not a
Hochschild class.

\paragraph{Relation to the other papers.} This note uses the definitions of \cite[\S\S2, 5, 10]{Conditionals} and of the
coupling 2-cells of \cite[\S7]{Attribution}. From the tower paper \cite{Tower} it uses only the context comonad, its distributive law and the stochastic change of base, in
\cref{sec:monad}. It does not use the obstruction tower, the empirical material of
\cite{ThreeLayers}, or the toric computations of \cite{Family}, except to name the variety in \cref{thm:ehrhart} and to compare
walls in \cref{prop:twobym}.

%=====================================================================
\section{The gluing algebra}\label{sec:gluing}
%=====================================================================

Let $f$ be a full-support distribution on $Y=\{1,\dots,n\}$, and $D=\diag(f)$.

\begin{definition}\label{def:Af}
$A_f=\{a\in\R^{n\times n}:a\mathbf 1=s\,f,\ \mathbf 1^\T a=s\,f^\T\text{ for some }s=s(a)\in\R\}$, with product
$a\cdot b=aD^{-1}b$ and unit $D$. The \emph{mass} is $s(a)$, and $I_f=\ker s$ is the space of \emph{circulations}.
\end{definition}

Couplings of $f$ with itself are the elements with $a\ge0$ and $s=1$, and their product is the gluing of
\cite[Thm.~2.3]{Conditionals}. Write $P_f=ff^\T$, the product coupling, and $e_f=D-P_f$.

\begin{theorem}[Structure of the gluing algebra]\label{thm:structure}
\begin{enumerate}[label=(\roman*)]
\item $P_f$ and $e_f$ are idempotents with $P_f+e_f=D$ and $P_f\cdot e_f=0$, and both are central.
\item $P_f\cdot a=a\cdot P_f=s(a)P_f$ and $e_f\cdot a=a-s(a)P_f$. Hence $A_f=\R P_f\oplus I_f$ as algebras, $I_f=e_fA_f$ is a
  unital algebra with unit $e_f$, and $s$ is the projection onto the first factor.
\item With $v=(\sqrt{f_y})_y$, the map $\varphi(a)=D^{-1/2}aD^{-1/2}$ is an isomorphism of $A_f$ onto
  $\{M:Mv=\lambda v,\ v^\T M=\lambda v^\T\}=\R\,vv^\T\oplus\operatorname{End}(v^\perp)$. So
  $A_f\cong\R\times M_{n-1}(\R)$, a semisimple algebra, and $\dim I_f=(n-1)^2$.
\end{enumerate}
\end{theorem}
\begin{proof}
(ii) $P_fD^{-1}a=f(\mathbf 1^\T a)=s(a)ff^\T$, and symmetrically. Then $e_f\cdot a=a-s(a)P_f$, which has mass $0$. (i) follows
from (ii) with $s(P_f)=1$ and $s(e_f)=0$, using $f^\T D^{-1}f=\sum_yf_y=1$. (iii) $\varphi(aD^{-1}b)=\varphi(a)\varphi(b)$ and
$\varphi(D)=1$. The conditions $a\mathbf 1=sf$ and $\mathbf 1^\T a=sf^\T$ become $\varphi(a)v=sv$ and $v^\T\varphi(a)=sv^\T$, since
$D^{1/2}\mathbf 1=v$ and $D^{-1/2}f=v$. If $Mv=\lambda v$ and $v^\T M=\lambda v^\T$, then $M$ maps $v^\perp$ into itself, so $M$ is
$\lambda vv^\T$ plus an endomorphism of $v^\perp$; and $\varphi(P_f)=vv^\T$.
\end{proof}

So the algebra is as simple as it could be, and none of the coupling geometry is visible in it: the couplings are the
nonnegative elements of mass one, and nonnegativity is not an algebraic condition here. The product coupling, the
Birch point of the independence model and the Kähler class of \cite[\S6]{Conditionals}, appears as the central idempotent
that projects onto the masses.

\paragraph{Several objects.} Let $\mathcal L$ be the linear category of \cite[Def.~10.1]{Conditionals}: objects are
full-support distributions on finite sets, $\mathcal L(f,g)$ consists of the matrices with row sums $sf$ and column sums
$sg$, and composition is $a\cdot b=a\,\diag(g)^{-1}b$.

\begin{proposition}[Splitting]\label{prop:split}
Every $a\in\mathcal L(f,g)$ is uniquely $a=s(a)\,fg^\T+x$ with $x$ a circulation, and
\[
\bigl(s\,fg^\T+x\bigr)\cdot\bigl(t\,gk^\T+y\bigr)=st\,fk^\T+x\cdot y .
\]
So each hom-space of $\mathcal L$ is the direct sum of $\R$ and the circulations, compatibly with composition: $\mathcal L$ is the
componentwise sum of the linearized indiscrete category, with every hom-space $\R$, and the category of circulations,
whose identity at $f$ is $e_f$.
\end{proposition}
\begin{proof}
$fg^\T\diag(g)^{-1}y=f(\mathbf 1^\T y)=0$ for a circulation $y$, and $x\,\diag(g)^{-1}gk^\T=(x\mathbf 1)k^\T=0$; the remaining term is
$st\,f(\mathbf 1^\T\diag(g)^{-1}g)k^\T=st\,fk^\T$.
\end{proof}

%=====================================================================
\section{The monad and the comonad}\label{sec:monad}
%=====================================================================

The gluing algebra is defined by matrices, but it comes from the distribution monad $\D$. This section says how, and
how far the context comonad of the tower paper \cite{Tower} reaches into it. Write $K^\#\colon\D(Y)\to\D(Y)$,
$K^\#(\rho)=K^\T\rho$, for the Kleisli extension of a kernel $K\colon Y\to\D(Y)$, written as a stochastic matrix. The Kleisli
extensions are exactly the morphisms of $\D$-algebras $\D(Y)\to\D(Y)$ between free algebras, that is, the affine self-maps of
the simplex. Let
\[
\operatorname{End}(f)=\{K\colon Y\to\D(Y)\ :\ K^\#f=f\},
\]
the monoid, under Kleisli composition, of $\D$-algebra endomorphisms of $\D(Y)$ that fix the point $f$.

\begin{proposition}[Couplings are invariant kernels]\label{prop:kleisli}
\begin{enumerate}[label=(\roman*)]
\item The map $\alpha\mapsto D^{-1}\alpha$ is an isomorphism of monoids from $\Cpl(f,f)$ under gluing onto $\operatorname{End}(f)$
  under Kleisli composition, with the unit $D$ going to the identity kernel. Its linear extension $a\mapsto D^{-1}a$ is an
  isomorphism of $A_f$ onto the algebra of matrices $M$ with $M\mathbf 1=s\mathbf 1$ and $M^\T f=sf$, and $A_f$ is the linear span of the
  couplings, so of $\operatorname{End}(f)$.
\item Let $\mathcal P$ be the category whose objects are the full-support probability spaces $(Y,f)$, that is, the objects of
  the coslice of the Kleisli category under the one-point set, restricted to full support, and whose morphisms are the
  kernels that push the first distribution to the second. The one-point space $(1,\delta)$ is a zero object of $\mathcal P$.
  Under (i), $\mathcal P(f,g)\cong\Cpl(f,g)$, and $\mathcal L$ is the category obtained from $\mathcal P$ by taking linear spans of the
  hom-sets.
\item The morphism of $\mathcal P$ that factors through the zero object is the constant kernel $y\mapsto g$, which corresponds to the
  product coupling $fg^\T$. So the first summand of \cref{prop:split} is spanned by the maps through the zero object, and
  $P_f$ is the endomorphism of $(Y,f)$ through it.
\item The mass is the pushforward along the unique morphism $(Y,f)\to(1,\delta)$: for this kernel $!$ the product action of
  \cref{sec:actions} gives $\pi_!(a)=\mathbf 1^\T a\mathbf 1=s(a)$, and it is multiplicative.
\end{enumerate}
\end{proposition}
\begin{proof}
(i) If $\alpha\in\Cpl(f,f)$ then $K=D^{-1}\alpha$ is nonnegative with $K\mathbf 1=D^{-1}f=\mathbf 1$ and $K^\T f=\alpha^\T\mathbf 1=f$; conversely
$\alpha=DK$ for $K\in\operatorname{End}(f)$ is a coupling. Gluing is $\alpha D^{-1}\beta=D(D^{-1}\alpha)(D^{-1}\beta)$, which is Kleisli composition of
the kernels in matrix form. The linear statement is the same computation without positivity. The couplings span $A_f$ because
$P_f$ has positive entries, so $P_f+tx$ is a coupling for every circulation $x$ and small $t$; the span therefore contains $\R P_f+I_f=A_f$.
(ii) A morphism $(1,\delta)\to(Y,f)$ is a kernel $1\to\D(Y)$ sending $\delta$ to $f$, so it is $f$ itself; a morphism $(Y,f)\to(1,\delta)$ is
the unique kernel to the one-point set. The isomorphism $\mathcal P(f,g)\cong\Cpl(f,g)$ is (i) with two objects, and composition in
$\mathcal L$ is the linear extension of composition in $\mathcal P$. (iii) The composite $Y\to1\to\D(Y)$ is $y\mapsto g$, with matrix $\mathbf 1g^\T$, and
$D\,\mathbf 1g^\T=fg^\T$. (iv) $\mathbf 1^\T aD^{-1}b\mathbf 1=s(a)\,f^\T D^{-1}f\,s(b)=s(a)s(b)$.
\end{proof}

So the gluing algebra is the linearized endomorphism monoid of a point of a free $\D$-algebra, and the structure theorem
reads categorically. The central idempotent $P_f$ is the endomorphism through the zero object, the mass is the
linearized map to the zero object, and the circulations are the complementary summand. The barycenter
$\mu_Y\colon\D^2(Y)\to\D(Y)$ plays no role in the splitting: the mass line comes from the zero object, not from the monad
multiplication.

\paragraph{The comonad.} The tower paper builds, for a finite base $B$, the context algebra $A_B=\D B\times\D^2B\times\D^3B\times\Sh(B)$
with the levelwise barycenter $\beta_B$, which makes it a product of free $\D$-algebras \cite[\S9]{Tower}. The comonad
$\C^B(Y)=Y\times A_B$ is the product (reader) comonad with parameter $A_B$, and the mixed distributive law
$\lambda^B\colon\D\C^B\Rightarrow\C^B\D$ of \cite[Thm.~9.8]{Tower} is $\lambda^B_Y(M)=(\D(\mathrm{pr}_1)M,\ \beta_B(\D(\mathrm{pr}_2)M))$. Mixed distributive
laws correspond to liftings of the comonad to the Eilenberg--Moore category \cite{PowerWatanabe2002}, and the lifting here sends a
$\D$-algebra $(X,\xi)$ to the product algebra $(X,\xi)\times(A_B,\beta_B)$. So the law is exactly the statement that $A_B$ is a
$\D$-algebra. The coalgebras of $\C^B$ are the sets over $A_B$, and the cofree coalgebra on $X$ is $X\times A_B$
\cite[Props.~13.1, 13.2]{Tower}. A kernel $k\colon B\to\D(B')$ acts by the $\D$-algebra map
$A_k=k^\#\times\D(k^\#)\times\D^2(k^\#)\times\Sh(k)$, with $A_{h\circ k}=A_hA_k$ \cite[Prop.~9.14]{Tower}.

\begin{proposition}[The comonad sees the monoid, not the algebra]\label{prop:comonad}
Take $B=Y$, and let $A^f_Y\subseteq A_Y$ be the contexts whose level-one component is $f$.
\begin{enumerate}[label=(\roman*)]
\item $K\mapsto A_K$ is an action of the monoid $\operatorname{End}(f)$ on $A^f_Y$, by $\D$-algebra maps. It is trivial on the level-one
  component.
\item On level one, $K\mapsto K^\#$ is affine in $K$. On level two it is not: for $K_1,K_2\in\operatorname{End}(f)$ with
  $\bar K=\frac12(K_1+K_2)$ and a context with level-two component $\Pi$, in general
  $\D(\bar K^\#)\Pi\ne\frac12\D(K_1^\#)\Pi+\frac12\D(K_2^\#)\Pi$.
\item Hence the action does not extend to a linear action of $A_f$ on the span of the contexts, compatible with
  \cref{prop:kleisli}(i). The algebra $A_f$ linearizes the monad side only; the comonad side sees the monoid
  $\operatorname{End}(f)$.
\end{enumerate}
\end{proposition}
\begin{proof}
(i) $A_K$ is a $\D$-algebra map and $A_{K_2\circ K_1}=A_{K_2}A_{K_1}$ by \cite[Prop.~9.14]{Tower}, and $K^\#f=f$. (ii) $K^\#\rho=K^\T\rho$ is linear in
$K$. For the second statement take $\Pi=\frac12\delta_{\rho_1}+\frac12\delta_{\rho_2}$ with $\frac12(\rho_1+\rho_2)=f$. Then
$\D(\bar K^\#)\Pi$ is supported on the two points $\bar K^\#\rho_i$, while the average of the $\D(K_j^\#)\Pi$ is supported on the four points
$K_j^\#\rho_i$, which are distinct for generic $K_1,K_2$. This is the phenomenon behind the failure of the shift law and of the level
separation in \cite[Thm.~8.3, Ex.~10.6]{Tower}: averaging inside a level is not averaging of levels. (iii) A linear action of the
span $A_f$ would restrict to an affine map on the convex set $\operatorname{End}(f)$, contradicting (ii).
\end{proof}

\begin{remark}[What the monadic reading does not give]\label{rem:monad-not}
Several stronger identifications suggest themselves, and each fails.
\begin{enumerate}[label=(\alph*),leftmargin=*]
\item \emph{The barycenter is not the mass projection.} $\mu_Y$ is a map $\D^2(Y)\to\D(Y)$ and does not act on $A_f$. The projection
  onto $\R P_f$ is \cref{prop:kleisli}(iv), induced by the map to the zero object.
\item \emph{$A_f$ is the span of $\operatorname{End}(f)$, not equal to it.} The monoid consists of the nonnegative elements of mass one;
  the algebra is its linear span, and it is in passing to the span that positivity, hence contextuality, is lost
  (\cref{cor:blind}).
\item \emph{No ideal filtration.} By \cref{thm:structure}, $A_f\cong\R\times M_{n-1}(\R)$ is a product of two simple algebras, so its
  two-sided ideals are $0$, $\R P_f$, $I_f$ and $A_f$; any single nonzero circulation generates $I_f$. There is no room for a
  three-step filtration by ideals, so the precision ladder of \cite[Thm.~3.5]{Attribution}, whose rungs are nested attribution sets, is not a filtration of $A_f$. The ladder lives in the 1-cell layer. A decomposition
  of $p$ with exact sources is a weight vector $(\lambda_1,\dots,\lambda_k,\mu)$ and an unmodeled law $U$ with
  $p=k^\#(\lambda,\mu)$ for the kernel $k\colon i\mapsto A_i$, $u\mapsto U$; so the attribution set is a union, over $U$, of fibres over
  $p$ of the Kleisli extension of a kernel, cut by the budget $\mu\le\bar\mu$ \cite[Def.~3.1, Thm.~3.2]{Attribution}. Sources known
  by their supports enlarge the set of admissible kernels. The rungs grow because more kernels are admitted and the budget
  grows, which is a statement about fibres of a 1-cell, not about the gluing of 2-cells.
\item \emph{The defect is not a Hochschild class.} The arity-three equation in \cref{thm:forced} is formally the Hochschild cocycle
  identity for $b_2$, but $I_{f'}$ is an $A_f$-bimodule through $F_1$ only when $F_1$ is multiplicative, and then $b_2=0$. In any case
  $A_f$ is separable, so its Hochschild cohomology vanishes in positive degrees for every bimodule, and no Hochschild class can
  carry the defect.
\item \emph{No explanation of non-unification.} The splitting is internal to the coupling layer. The non-unification results of
  \cite{ThreeLayers} concern different coefficient systems, and nothing here maps them to $A_f$.
\end{enumerate}
\end{remark}

%=====================================================================
\section{How kernels act}\label{sec:actions}
%=====================================================================

A kernel $h\colon Y\to\D(Z)$, written as a stochastic matrix, acts on matrices in two ways \cite[Def.~10.1]{Conditionals}.
The \emph{product action} $\pi_h(a)=h^\T ah$ runs $h$ independently on both coordinates. The \emph{whiskering} $w_h$ copies
$h$ on equal pairs and runs it independently on unequal pairs, which is the postcomposition of the witness layer
\cite[Thm.~7.6]{Attribution}. Both preserve mass and push margins forward.

\begin{proposition}[The actions in split coordinates]\label{prop:actions}
Write $h_y$ for the row of $h$ at $y$.
\begin{enumerate}[label=(\roman*)]
\item $\pi_h(fg^\T)=(h^\T f)(h^\T g)^\T$, and $\pi_h$ maps circulations to circulations. So $\pi_h$ respects the splitting.
\item For every $a$, $w_h(a)=\pi_h(a)+\sum_ya_{yy}\bigl(\diag(h_y)-h_yh_y^\T\bigr)$, and each term of the sum is a circulation.
  In particular $w_h(P_f)=P_{h^\T f}+c_h(f)$ with the \emph{copy correction} $c_h(f)=\sum_yf_y^2\bigl(\diag(h_y)-h_yh_y^\T\bigr)$.
\item The multiplicativity defect of the whiskering on the product coupling is
  $b_2(P_f,P_f)=w_h(P_f\cdot P_f)-w_h(P_f)\cdot w_h(P_f)=c_h(f)-c_h(f)\cdot c_h(f)$.
\end{enumerate}
\end{proposition}
\begin{proof}
(i) is direct, and $\pi_h$ of a matrix with zero row and column sums has zero row and column sums because $h$ is stochastic.
(ii) The two actions agree off the diagonal of $a$; on an equal pair $(y,y)$ copying replaces $a_{yy}h_yh_y^\T$ by
$a_{yy}\diag(h_y)$, and $\diag(h_y)-h_yh_y^\T$ has zero row and column sums. The diagonal of $P_f$ is $(f_y^2)_y$. (iii) $P_f\cdot P_f=P_f$, and $P_{h^\T f}$ annihilates circulations by
\cref{thm:structure}(ii).
\end{proof}

So the whiskering mixes the two factors of \cref{thm:structure}, and the copy correction, a weighted sum of the covariance
matrices of the rows of $h$, measures by how much. The product action respects the factors but does not preserve the unit,
since $\pi_h(D)=h^\T Dh\ne\diag(h^\T f)$ unless $h$ is deterministic. This is the trade-off of \cite[Thm.~10.2]{Conditionals}
seen in coordinates.

%=====================================================================
\section{The square-zero cone}\label{sec:cone}
%=====================================================================

Let $f'$ be a distribution and $C=A_{f'}\oplus\varepsilon I_{f'}$, with $A_{f'}$ in degree $0$ and $\varepsilon I_{f'}$ in degree $-1$. The product is
$a\cdot\varepsilon x=\varepsilon(a\cdot x)$, $\varepsilon x\cdot a=\varepsilon(x\cdot a)$ and $\varepsilon x\cdot\varepsilon y=0$, which is associative and unital because $I_{f'}$ is a
two-sided ideal, and the differential is $m_1(\varepsilon x)=x$, which satisfies the Leibniz rule. So $C$ is a differential graded
algebra, whose underlying complex is the mapping cone of $I_{f'}\hookrightarrow A_{f'}$. We use the sign conventions of \cite{Keller2001}.

\begin{theorem}[The functor is forced]\label{thm:forced}
Let $F_1\colon A_f\to A_{f'}$ be linear with $s(F_1a)=s(a)$. Then there is exactly one $A_\infty$-functor $F\colon A_f\to C$ with
first component $F_1$. Its components are $F_2=\varepsilon\,b_2$, with $b_2(a,b)=F_1(a\cdot b)-F_1(a)\cdot F_1(b)$, and $F_n=0$ for
$n\ge3$.
\end{theorem}
\begin{proof}
The component $F_n$ has degree $1-n$. For $n\ge3$ the cone has nothing in that degree, so $F_n=0$. The arity-two equation
reads $m_1F_2(a,b)=b_2(a,b)$. Since $F_1$ preserves mass, $b_2(a,b)$ has mass $s(a)s(b)-s(a)s(b)=0$, so it lies in $I_{f'}$.
Since $F_2$ has degree $-1$ it lies in $\varepsilon I_{f'}$, where $m_1$ is injective, and so $F_2=\varepsilon b_2$ is the only solution. The arity-three
equation is $F_1a\cdot b_2(b,c)-b_2(ab,c)+b_2(a,bc)-b_2(a,b)\cdot F_1c=0$. Expanding $b_2$, the terms $F_1(abc)$ and
$F_1a\,F_1b\,F_1c$ cancel, using only linearity of $F_1$ and associativity. In arity four the only surviving term is a
product $F_2\cdot F_2\in\varepsilon^2I_{f'}=0$, and in higher arity every term contains some $F_k$ with $k\ge3$.
\end{proof}

Whiskering and the product action are two instances, and so are their convex combinations. The cone is therefore not
specific to whiskering: it is the receptacle of every mass-preserving linear map, which it turns into an $A_\infty$-functor
by recording the multiplicativity defect.

\begin{theorem}[Homotopy content]\label{thm:homotopy}
\begin{enumerate}[label=(\roman*)]
\item $H^0(C)=A_{f'}/I_{f'}\cong\R$ and $H^{-1}(C)=0$, so the map $q(a+\varepsilon x)=s(a)$ is a quasi-isomorphism $C\to\R$ of
  differential graded algebras. In particular the minimal model of $C$ is $\R$.
\item For the functor $F$ of \cref{thm:forced}, $q\circ F$ is the mass $s\colon A_f\to\R$, a strict algebra map.
\item Let $G\colon A_f\to E$ be any $A_\infty$-functor into an $A_\infty$-algebra $E$ such that $G_1(e_f)$ is a boundary. Then
  $G_1(x)$ is a boundary for every circulation $x$, and $H(G_1)$ factors through the mass.
\end{enumerate}
\end{theorem}
\begin{proof}
(i) $m_1$ maps $\varepsilon I_{f'}$ isomorphically onto $I_{f'}$, and $A_{f'}/I_{f'}\cong\R$ by \cref{thm:structure}. The map $q$ is
multiplicative because $s$ is and because $q$ vanishes on $\varepsilon I_{f'}$. A dg algebra with cohomology $\R$ in degree $0$ has
minimal model $\R$. (ii) $q(F_1a)=s(F_1a)=s(a)$ and $q(F_2)=0$. (iii) An $A_\infty$-functor induces an algebra map on homology.
For a circulation $x$, $x=e_f\cdot x$ by \cref{thm:structure}(ii), so $H(G_1)(x)=H(G_1)(e_f)\,H(G_1)(x)=0$.
\end{proof}

\begin{corollary}\label{cor:blind}
Let $G\colon A_f\to E$ be an $A_\infty$-functor such that $G_1(e_f)$ is a boundary. Then $G$ is homotopic, after passing to a
minimal model of $E$, to a strict algebra map that factors through the mass $s$. So every homotopy invariant of $G$ is a
function of masses alone.
\end{corollary}
\begin{proof}
By \cref{thm:structure}, $A_f$ is a product of matrix algebras over $\R$, hence separable, so its Hochschild cohomology with
coefficients in any bimodule vanishes in positive degrees. The obstructions to removing the higher components of an
$A_\infty$-functor out of $A_f$, after passing to a minimal model of the target, lie in these groups, so $G$ is homotopic to a
strict map. That map is $H(G_1)$, which factors through $s$ by \cref{thm:homotopy}(iii).
\end{proof}

In particular no such invariant detects contextuality, which is an obstruction to positivity: signed measures always glue
\cite[Prop.~35.1]{ThreeLayers}, and the parity horns of \cite[Thm.~4.3]{Conditionals} are compatible families with contextual
fraction $1-2^{2-n}$ that a linear invariant cannot distinguish from noncontextual ones.

So the coupling information of the $A_\infty$-functor lives at chain level, in $F_2=\varepsilon b_2$, and nowhere in its homotopy
type. This settles \cite[Question~9.4(i)]{Conditionals} in the negative for homotopy-invariant information. Whether the free
cone carries chain-level higher components given by the higher Boolean cumulants remains open (\cref{sec:open}).

%=====================================================================
\section{Couplings with kernels do not form a double category}\label{sec:double}
%=====================================================================

Take full-support probability spaces $(Y,p)$ as objects, couplings $\alpha\in\Cpl(p,q)$ as horizontal arrows composed by gluing, and
kernels $k$ with $k^\T p=p'$ as vertical arrows. A square with horizontal arrows $\alpha\in\Cpl(p,q)$ and $\alpha'\in\Cpl(p',q')$ and
vertical kernels $k,l$ exists when $\alpha'=k^\T\alpha l$; squares are a property, so the structure would be a thin double category.
Squares compose vertically, because $\pi$ is functorial in the kernel. For them to compose horizontally, gluing must commute
with pushforward along the middle kernel.

\begin{proposition}[Interchange]\label{prop:double}
Let $\alpha\in\Cpl(p,q)$, $\beta\in\Cpl(q,r)$, and kernels $k,l,m$ with $l$ acting on the middle space. The identity
\[
k^\T(\alpha\cdot\beta)m=(k^\T\alpha l)\cdot(l^\T\beta m),
\]
where the right side glues over $l^\T q$, holds whenever $l$ is a bijection, and more generally whenever $l$ is a function such that on each of its fibres the
conditional law of the first coordinate given the middle one, or that of the last coordinate given the middle one, is
constant. It fails for some constant kernels $l$ and for some non-injective functions $l$, so squares are
not closed under horizontal composition. The difference of the two sides is always a circulation.
\end{proposition}
\begin{proof}
Since $k$ and $m$ act on the outer coordinates, it suffices to take them to be identities. For a function $l$ write $y'$ for a
fibre and $q(y')$ for its mass. With $\alpha_{\cdot y}$ the columns of $\alpha$ and $\beta_{y\cdot}$ the rows of $\beta$, the two sides are
\[
\sum_{y'}\sum_{y\in y'}\frac{\alpha_{\cdot y}\beta_{y\cdot}}{q_y}
\quad\text{and}\quad
\sum_{y'}\frac1{q(y')}\Bigl(\sum_{y\in y'}\alpha_{\cdot y}\Bigr)\Bigl(\sum_{y\in y'}\beta_{y\cdot}\Bigr).
\]
If $\alpha_{\cdot y}=q_y c_{y'}$ for $y\in y'$, both fibre terms equal
$c_{y'}\sum_{y\in y'}\beta_{y\cdot}$, and symmetrically for $\beta$. A bijection has singleton fibres. For a counterexample with a function, let
$p=q=r$ be uniform on $\{0,1\}$, let $\alpha=\beta$ be the diagonal coupling, let $k,m$ be identities and let $l$ collapse
$\{0,1\}$ to a point. The left side is the diagonal coupling, since gluing two copies of the identity gives the identity. On
the right, $k^\T\alpha l$ and $l^\T\beta m$ are product couplings with a one-point middle space, so their gluing is the product
coupling $\frac14\mathbf 1\mathbf 1^\T$. For a genuinely stochastic kernel, take instead $l$ with both rows uniform on $\{0,1\}$: then
$\alpha l$ and $l^\T\beta$ are again product couplings, so the right side is still $\frac14\mathbf 1\mathbf 1^\T$. Both sides have the margins
$k^\T p$ and $m^\T r$, so their difference is a circulation.
\end{proof}

So the structure is a double category, strict and thin, when the vertical arrows are restricted to bijections. The
sufficient condition says that interchange holds when the middle map merges only values that are interchangeable for one
of the outer variables, that is, when $l(Y)$ is sufficient for $Y$ with respect to $X$ or to $Z$. With all
kernels, horizontal composition of squares is not defined in general, and since squares are a property there are no
comparison cells that could make it lax. The defect is a circulation of the same kind as the whiskering defect, so a
double-category formulation does not remove the obstruction of \cite[Thm.~7.6(v)]{Attribution}.

%=====================================================================
\section{The ring of contingency tables}\label{sec:tables}
%=====================================================================

The gluing algebra forgets the geometry of couplings. A different algebra keeps it. Let $f,g$ be rational margins on
$\{1,\dots,n\}$ and $\{1,\dots,m\}$ with common denominator $N$, so that $Nf$ and $Ng$ are integer vectors with the same total.
The polytope $P=N\cdot\Cpl(f,g)$ consists of the nonnegative real $n\times m$ tables with row sums $Nf$ and column sums $Ng$.
Its lattice is the set of integer tables, and the lattice points of $kP$ are the integer tables with margins $kNf$ and
$kNg$.

\begin{theorem}[The ring of contingency tables]\label{thm:ehrhart}
\begin{enumerate}[label=(\roman*)]
\item $P$ is a lattice polytope of dimension $(n-1)(m-1)$ when $f,g$ have full support, and it has the integer
  decomposition property: every integer table in $kP$ is a sum of $k$ integer tables in $P$.
\item The Ehrhart ring $S_P=\bigoplus_{k\ge0}\R[kP\cap\Z^{n\times m}]$, graded by $k$ with multiplication given by adding tables, is
  generated in degree one. Its Hilbert function $k\mapsto\dim S_{P,k}$ is the Ehrhart polynomial of $P$, of degree
  $(n-1)(m-1)$, with constant term $1$ and leading coefficient the volume of $P$ normalized to the lattice of integer
  circulations.
\item $\operatorname{Proj}S_P$ is the projective toric variety of $P$, embedded projectively normally. For nondegenerate
  margins it is the smooth toric variety $X_{f,g}$ of \cite[Thm.~6.4]{Conditionals}, and $S_P$ is the ring of sections of the powers
  of the ample line bundle of $P$, whose class in $\operatorname{Pic}(X_{f,g})\otimes\mathbb Q$ is $N$ times the class of the product coupling.
\item The degree-one basis $P\cap\Z^{n\times m}$ is the set of integer tables with margins $Nf$ and $Ng$, the sample space of
  Fisher's exact test, on which the conditional law of a table given its margins under independence is the multivariate
  hypergeometric law, proportional to $1/\prod_{ij}t_{ij}!$. The basic moves $\delta_{ij}+\delta_{i'j'}-\delta_{ij'}-\delta_{i'j}$
  connect this set, and each is an integer circulation.
\end{enumerate}
\end{theorem}
\begin{proof}
(i) The constraint matrix of a transportation polytope is the incidence matrix of a complete bipartite graph, which is
totally unimodular, so the vertices of $P$ are integral \cite{Schrijver1986}, and polyhedra defined by totally unimodular
systems with integral right-hand sides have the integer decomposition property \cite{BaumTrotter1978}. The dimension is that of
the circulation space. (ii) Generation in degree one is (i). The Hilbert function counts lattice points of $kP$, and the
properties of the Ehrhart polynomial are standard \cite{BeckRobins2015}. (iii) A lattice polytope with the integer
decomposition property is normal, hence very ample, and $\operatorname{Proj}$ of its Ehrhart ring is the toric variety of its
normal fan, embedded projectively normally, and the sections of the $k$-th power of its line bundle have the lattice
points of $kP$ as a basis \cite[Ch.~2 and Prop.~4.3.3]{CLS2011}. The normal fan of $P$ is that of $\Cpl(f,g)$, which for
nondegenerate margins is the smooth fan of $X_{f,g}$. The support numbers of $P$ are $N$ times the entries of the product
coupling, which give the class of $\Cpl(f,g)$ \cite[Prop.~8.1]{Conditionals}; an integral representative comes from any
lattice point of $P$. (iv) The hypergeometric law is classical \cite{Agresti2013}. For two-way tables the basic moves
form a Markov basis \cite{DiaconisSturmfels1998}; each has zero row and column sums.
\end{proof}

\begin{example}\label{ex:hexagon}
For $f=(\frac12,\frac12)$ and $g=(\frac13,\frac13,\frac13)$, take $N=6$. The tables with margins $(3,3)$ and $(2,2,2)$ are
$7$ lattice points of a hexagon. The numbers of tables with margins $k(3,3)$ and $k(2,2,2)$ are $7,19,37,61,91$ for $k=1,\dots,5$,
equal to $3k^2+3k+1$. The leading coefficient $3$ is the lattice area of the hexagon, which is $N^2=36$ times the area
$\frac1{12}$ of $\Cpl(f,g)$ in \cite[\S8]{Conditionals}. The hexagon has one interior point, so it is the reflexive hexagon of the
del Pezzo surface of degree $6$, and its $h^*$-vector, read off from $\sum_k L(k)t^k=h^*(t)/(1-t)^3$, is $(1,4,1)$, of sum $6$, the
normalized area. For margins $(5,7)$ and $(3,4,5)$ the counts give $10k^2+6k+1$, the area of that polygon is $10$, and
$h^*=(1,14,5)$.
\end{example}

For two rows the dependence on the margins can be written down completely. Write $r=(r_1,r_2)$ and $c=(c_1,\dots,c_m)$ for
integer margins with $r_1+r_2=\sum_jc_j$, $N(r,c)$ for the number of integer tables, and $c(S)=\sum_{j\in S}c_j$. The walls of
\cite[Def.~6.3]{Conditionals} for two rows are the hyperplanes $r_1=c(S)$ with $\emptyset\ne S\subsetneq\{1,\dots,m\}$, since $r_2=c(S)$
is $r_1=c(S^c)$.

\begin{proposition}[Two-row tables across the walls]\label{prop:twobym}
\begin{enumerate}[label=(\roman*)]
\item $N(r,c)=\sum_{S}(-1)^{|S|}\binom{r_1-c(S)-|S|+m-1}{m-1}$, the sum over $S=\emptyset$ and the subsets
  $S\subseteq\{1,\dots,m\}$ with $c(S)<r_1$.
\item On each chamber $N$ is a polynomial in $(r_1,c)$ of degree $m-1$. Crossing the wall $r_1=c(S)$ in the direction of
  increasing $r_1$ adds the polynomial $(-1)^{|S|}\binom{r_1-c(S)-|S|+m-1}{m-1}$, which vanishes on the lattice points with
  $|S|-m+1\le r_1-c(S)\le|S|-1$. So the polynomials of adjacent chambers agree on a strip along the wall.
\item For $m=3$ and $|S|=1$, the polygon on the far side is obtained by cutting a corner of lattice depth $d=r_1-c_j$, and the
  count drops by $\binom{d+1}2$. This is the Riemann--Roch identity $\chi(\pi^*L-dE)=\chi(L)-\binom{d+1}2$ for the toric blow-up
  $\pi$ at the corresponding fixed point with exceptional curve $E$, and its leading term $\frac12d^2$ is, after rescaling by $N$, the area
  jump $\frac12s^2$ of \cite[Prop.~10.6(i)]{Conditionals}.
\end{enumerate}
\end{proposition}
\begin{proof}
(i) A table is determined by its first row $t$, with $0\le t_j\le c_j$ and $\sum t_j=r_1$. Without the upper bounds there are
$\binom{r_1+m-1}{m-1}$ such rows, and inclusion--exclusion over the set $S$ of violated bounds, substituting $t_j=c_j+1+t'_j$ for
$j\in S$, gives the terms with $r_1-c(S)-|S|\ge0$. A term with $c(S)<r_1$ but $r_1-c(S)-|S|<0$ has top argument between $0$ and $m-2$,
so it is zero, and including it changes nothing. (ii) On a chamber the set of subsets with $c(S)<r_1$ is constant, so (i) is a polynomial.
Crossing $r_1=c(S)$ adds exactly the term for $S$, and $\binom{x+m-1}{m-1}$ vanishes for $x=-1,\dots,-(m-1)$. (iii) For $m=3$ the
polygon is the triangle $\{t\ge0,\ t_1+t_2+t_3=r_1\}$ cut by the bounds $t_j\le c_j$, and the term for $S=\{j\}$ is minus the number
$\binom{d+1}2$ of lattice points with $t_j>c_j$, a corner triangle. For the toric surface, $(\pi^*L-dE)^2=L^2-d^2$ and
$(\pi^*L-dE)\cdot K=L\cdot K+d$, since $K=\pi^*K+E$ and $E^2=-1$; Riemann--Roch gives the identity, and the lattice points of the
polygon count $\chi$ because the line bundle is nef. The leading term of $\binom{d+1}2$ is $\frac12d^2$.
\end{proof}

General transportation polytopes behave the same way in outline: the number of tables is a vector partition function,
polynomial on the chambers because the constraint matrix is unimodular \cite{Sturmfels1995}, and its jumps across walls are
given by the wall-crossing formula of Paradan as made explicit by Boysal and Vergne \cite{BoysalVergne2009}.

\begin{example}[Fisher's exact test on synthetic data]\label{ex:fisher}
The first $40$ items of the synthetic moderation data of \cite[\S12]{Conditionals} give the label-by-action table
\[
T=\begin{pmatrix}10&8&1\\2&5&14\end{pmatrix},\qquad\text{margins }(19,21)\text{ and }(12,13,15).
\]
Its fibre, the degree-one basis of the ring of $\frac1{40}\cdot$ these margins, has $151$ tables. The hypergeometric probabilities sum
to $1$, and the exact $p$-value, the total probability of tables no more likely than $T$, is $9.8\times10^{-5}$. The
$\chi^2$ approximation with $2$ degrees of freedom gives $1.8\times10^{-4}$.
\end{example}

The three appearances of the circulations now line up. They are the ideal $I_f$ of the gluing algebra (\cref{sec:gluing}),
the tangent space of the coupling polytope, whose dimension $(n-1)(m-1)$ is both the degree of the Ehrhart polynomial and
the number of degrees of freedom of the test of independence \cite[Prop.~12.1]{Conditionals}, and the lattice of Markov moves
that connects the sample space of the exact test. The first appearance is homotopically trivial (\cref{thm:homotopy}). The
other two carry the geometry, because they keep positivity: a table must stay nonnegative under a move, and the polytope
is cut out by nonnegativity.

%=====================================================================
\section{Open problems}\label{sec:open}
%=====================================================================

\begin{enumerate}[leftmargin=*]
\item \emph{The free cone at chain level.} In the free cone of $I_f$, products of homotopies survive, and the functor of
  \cref{thm:forced} acquires arity-four terms. Are the higher components given by the higher Boolean cumulants of
  whiskering? By \cref{thm:homotopy} any answer is invisible up to homotopy.
\item \emph{A positive homotopy theory.} \Cref{cor:blind} locates contextuality in positivity: every linear invariant is blind
  to it, and every horn fills with signed measures \cite[Prop.~35.1]{ThreeLayers}. Find a homotopy theory of cones of
  measures, for example through simplicial distributions, in which unfillable horns of the coupling nerve
  \cite[\S4]{Conditionals} are nontrivial classes. The concrete test case is the family of parity horns of
  \cite[Thm.~4.3]{Conditionals}, unfillable with contextual fraction $1-2^{2-n}$, whose secondary obstruction in the tower is
  $o_3^-=2\,\mathrm{CF}$ \cite[Thm.~11.2]{Conditionals}. A theory that sees them would be a positivity-preserving refinement of the
  square-zero cone: it must agree with the cone on signed measures, where it can only see masses, and it should recover the
  contextual fraction, or $o_3^-$, as the size of a class.
\item \emph{Ehrhart data across walls.} \Cref{prop:twobym} settles two rows. For $n\times m$ tables, compute the chamber
  polynomials and the $h^*$-vectors of $N\cdot\Cpl(f,g)$ chamber by chamber, and match the wall-crossing terms of
  \cite{BoysalVergne2009} with the variation of GIT and the central-charge jumps of \cite[\S10]{Conditionals} and
  \cite[\S3]{Family}, as \cref{prop:twobym}(iii) does for one blow-up. This is the algebraic-statistical counterpart of the
  wall-crossing of the companion paper: the number of tables with given margins is the Euler characteristic of a line
  bundle, and its jumps should be those of the derived categories across the wall.
\item \emph{Sections and potentials.} On $X_{f,g}$ the degree-one sections are indexed by tables. Relate the exact test to
  the disc potential of \cite[\S7]{Conditionals}, for example through the leximin coupling, which lies over the
  Floer-nontrivial fibres at the sampled points of \cite[Obs.~10.10]{Conditionals}.
\end{enumerate}

%=====================================================================
\section{Computational checks}\label{sec:code}
%=====================================================================

\texttt{check\_coupling\_algebra.py} runs fourteen checks, all passing:
\begin{itemize}[leftmargin=*,nosep]
\item A1--A2: \cref{thm:structure} for $n=2,\dots,5$, and \cref{prop:split} on objects of sizes $2,3,4$;
\item A3: \cref{prop:actions};
\item A4: \cref{thm:forced} for convex combinations of the two actions;
\item A5: \cref{thm:homotopy};
\item A6: \cref{prop:double};
\item A7--A9: \cref{thm:ehrhart}(i), (ii) and (iv) on two $2\times3$ examples and a $3\times3$ example;
\item A10: \cref{ex:fisher}, in floating point;
\item A11: \cref{prop:kleisli} for $n=3,4$, including the dimension of the span;
\item A12: \cref{prop:comonad}, with an explicit level-two context;
\item A13: \cref{rem:monad-not}(c), the sufficient condition of \cref{prop:double} on both sides, and its failure for a generic merge;
\item A14: \cref{prop:twobym} against enumeration for $m=3,4$, the chamber polynomials and corner cuts, and the $h^*$-vectors of
  \cref{ex:hexagon}.
\end{itemize}
Everything except A10 uses exact rational arithmetic.

\begin{thebibliography}{99}

\bibitem{Agresti2013}
A.~Agresti, \emph{Categorical Data Analysis}, 3rd ed., Wiley, 2013.

\bibitem{Attribution}
V.~K. Awasthi, Precise attribution of model interventions: attribution sets, mechanism reach, path defects and gluing.
Technical report, Johns Hopkins University, October 2026. Draft 5.

\bibitem{Conditionals}
V.~K. Awasthi, Conditionals, couplings and toric geometry: kernels, the coupling nerve, cumulants and a {F}ukaya
programme. Technical report, Johns Hopkins University, October 2026. Revised draft 4.

\bibitem{Family}
V.~K. Awasthi, Coupling polytopes and toric geometry: the $2\times3$ family, the square case, higher levels and stability.
Companion paper, Johns Hopkins University, October 2026.

\bibitem{ThreeLayers}
V.~K. Awasthi, Three layers of realisation: hierarchy recovery, holonomy, information, contextuality and categorical
dynamics. Technical report, Johns Hopkins University, revised October 2026. DOI: 10.13140/RG.2.2.24274.52163.

\bibitem{Tower}
V.~K. Awasthi, The three-level posterior tower: hierarchical {B}ayesian inference and mixed distributive laws.
Technical report, Johns Hopkins University, October 2026. Revised version 10.

\bibitem{BaumTrotter1978}
S.~Baum and L.~E. Trotter, Integer rounding and polyhedral decomposition for totally unimodular systems. In
\emph{Optimization and Operations Research}, Lecture Notes in Economics and Mathematical Systems 157, Springer, 1978,
15--23.

\bibitem{BeckRobins2015}
M.~Beck and S.~Robins, \emph{Computing the Continuous Discretely}, 2nd ed., Springer, 2015.

\bibitem{BoysalVergne2009}
A.~Boysal and M.~Vergne, Paradan's wall crossing formula for partition functions and Khovanski--Pukhlikov differential
operator. \emph{Ann. Inst. Fourier} \textbf{59} (2009), 1715--1752.

\bibitem{CLS2011}
D.~A. Cox, J.~B. Little and H.~K. Schenck, \emph{Toric Varieties}, Graduate Studies in Mathematics 124, American
Mathematical Society, 2011.

\bibitem{DiaconisSturmfels1998}
P.~Diaconis and B.~Sturmfels, Algebraic algorithms for sampling from conditional distributions. \emph{Ann. Statist.}
\textbf{26} (1998), 363--397.

\bibitem{Keller2001}
B.~Keller, Introduction to $A$-infinity algebras and modules. \emph{Homology Homotopy Appl.} \textbf{3} (2001), 1--35.

\bibitem{PowerWatanabe2002}
J.~Power and H.~Watanabe, Combining a monad and a comonad. \emph{Theoret. Comput. Sci.} \textbf{280} (2002), 137--162.

\bibitem{Schrijver1986}
A.~Schrijver, \emph{Theory of Linear and Integer Programming}, Wiley, 1986.

\bibitem{Sturmfels1995}
B.~Sturmfels, On vector partition functions. \emph{J. Combin. Theory Ser. A} \textbf{72} (1995), 302--309.

\end{thebibliography}
\end{document}
